



%Below is a summary of the results from the experiments described in \S\ref{sec:5-chapter-experimental-setup}. 



\section{Results}\label{sec:5-chapter-results}


    

\subsection[Current PLS evaluation standards]{Current PLS evaluation standards \hyperref[rq1]{\rqone}}\label{subsec:results-rqone}

\begin{figure}
\centering
    \includegraphics[width=0.95\linewidth]{research-chapters/5-chapter/figures/count_acl_eval.pdf}
    \caption[Evaluation metrics used by papers published in the ACL Anthology.]{Evaluation metrics used by papers published in the ACL Anthology. We report the count of papers using each method out of a total of 18 papers. We additionally report the evaluation strategies used by PLS shared tasks and the number of participants.}
    \label{fig:lit-survey-results}
\end{figure}
We found 18 ACL Anthology papers on the task of PLS and 3 shared tasks, representing 81 additional papers. \Cref{fig:lit-survey-results} shows the literature survey results, excluding metrics used by a single paper. FKGL is the most popular metric, followed by CLI and DCRS. LM-based evaluations are uncommon (4 of the 18 papers). The shared task BioLaySumm used FKGL and DCRS for both years, adding CLI in 2024. BioLaySumm 2025 is ongoing at the time of writing; the organizers plan to use FKGL, DCRS, and CLI. Our survey shows that PLS is an increasingly popular topic of study, as the number of participants in shared tasks increased from 8 in 2020 to 53 in 2024, emphasizing the importance of PLS evaluation. Less popular metrics include GFI, ARI, lexical proxies (e.g., number of sentences in a document), and FRE. In \cref{subsec:results-rqtwo}, we place the highest importance on the results of the most commonly used evaluation metrics: FKGL, CLI, and DCRS.






\subsection[Comparing traditional readability metrics to human judgments]{Comparing traditional readability metrics to human judgments \hyperref[rq2]{\rqtwo}}\label{subsec:results-rqtwo}

In \Cref{subtab:trad-corr-results}, we report the Pearson and Kendall-Tau correlation of 8 traditional readability metrics with human judgments. We find that 6 of the 8 metrics have less than $0.3$ Pearson correlation with human judgments. DCRS and CLI have the highest correlation, achieving $0.2$ Pearson points higher correlation than the most popular metric, FKGL (\Cref{subsec:results-rqone}). FKGL receives only $0.16$ Pearson and $0.08$ Kendall-Tau correlation, indicating little to no correlation with human judgment. 

\Cref{tab:example-summ} shows an example summary and readability scores, along with its human judgment. The human annotators gave the example summary an average rating of 4.05/5; they found the text fairly readable. However, the majority of traditional metrics give the summary poor readability scores: college level or higher. This is likely because the text includes domain-specific vocabulary, such as ``acute respiratory distress syndrome (ARDS),'' which is penalized by traditional metrics. 
Traditional readability metrics do not account for elements of the summary that make it more readable, such as defining ARDS as ``a very serious lung disease'' and explaining the scientists' motivation to ``test a new method of lung damage diagnosis.''

\begin{table}[!t]
\centering
    \begin{subtable}{.4\textwidth}
        \centering \footnotesize
        \begin{tabular}{l|ll}
            \textbf{Metric} & \textbf{Pearson} & \textbf{Kendall Tau} \\ \hline
             FKGL &  0.16 & 0.08 \\
            CLI & 0.36 & 0.20 \\
            \textbf{DCRS} & \textbf{0.37} & \textbf{0.24} \\
            GFI & 0.21 & 0.11 \\
            ARI & 0.10 & 0.02 \\
            FRE & 0.29 & 0.15 \\
            Spache & 0.13 & 0.04 \\
            LW & -0.06 & -0.03 \\
        \end{tabular}
        \caption{Traditional metric scores correlation with human judgment.}
        \label{subtab:trad-corr-results}
    \end{subtable}    
    \hspace{4mm}
    \begin{subtable}{.4\textwidth}
        \centering \footnotesize
            \begin{tabular}{l|ll}
            \textbf{Model} & \textbf{Pearson} & \textbf{Kendall Tau} \\ \hline
            % OLMo 2 7B & 0.15 & 0.10 \\
            Mistral 7B & 0.52 & 0.44 \\
            Mixtral 7B & 0.54 & 0.41 \\
            Gemma 7B & 0.54 & 0.43 \\
            Llama 3.1 8B & 0.45 & 0.34 \\
            \textbf{Llama 3.3 70B} & \textbf{0.56} & \textbf{0.35} \\
            \multicolumn{3}{l}{} \\
            \multicolumn{3}{l}{} \\
            \multicolumn{3}{l}{} \\
        \end{tabular}
        \caption{LM scores correlation with human judgment.}\label{subtab:lm-corr-results}
    \end{subtable} 
    % \vspace{-7mm}
    \caption[Correlation of each readability metric with human judgment]{We report the Pearson and Kendall-Tau correlation of each metric with human judgment. Tab.\ref{subtab:trad-corr-results} contains the correlation of traditional readability metrics with human judgment. DCRS and CLI have the highest correlation with human judgment. Notably, the most popular metric, FKGL, as shown in \S\ref{subsec:results-rqone}, has low correlation with human judgment. Tab.\ref{subtab:lm-corr-results} contains the correlation of LM models as evaluators with human judgment. All 5 models achieve higher correlation than all of the traditional metrics.
    }
    \label{tab:corr-results}

\end{table}


\begin{table*}[!htb]
\centering
    \begin{subtable}{.95\textwidth}
      \centering 
        \begin{tabular}{p{\textwidth}}
            \textexample{Scientists create a device which can detect the onset of acute respiratory distress syndrome (ARDS), a very serious lung disease, by measuring chemicals in patients’ exhaled breath The researchers wanted to test a new method of lung damage diagnosis by analyzing patient breath samples. In particular, the researchers were looking for better ways to detect acute respiratory distress syndrome (ARDS), a form of lung injury that causes inflammation and severe damage. [...]  a much larger group of test subjects is necessary to further validate their method. This new method of breath analysis could be a noninvasive, cost effective way to diagnose and track ARDS, and could potentially be modified to screen for other serious conditions as well.}
        \end{tabular}
        \caption{Excerpt of an example summary. This summary is written by an expert and is labeled as a low complexity summary.}\label{subtab:ex-summ}
    \end{subtable}%
    \vspace{4mm}
    \quad
    \begin{subtable}[t]{.65\textwidth}
        \centering \footnotesize
        \begin{tabular}{r|lll}
        Metric & Score & S-12 & US Grade Level \\ \hline
        FKGL $\downarrow$ & 13.9 & 12 & College \\
        CLI $\downarrow$ & 12.7 & 12 & College \\
        DCRS $\downarrow$ & 11.3 & 8.9 & College graduate \\
        GFI $\downarrow$ & 18.6 & 12 & Above college graduate \\
        ARI $\downarrow$ & 16.7 & 13 & College graduate \\
        FRE $\uparrow$ & 50.2 & 50 & 12th grade \\
        Spache $\downarrow$ & 8.7 & 12 & 9th grade \\
        LW $\uparrow$ & 19.5 & 60 & College graduate \\

        \end{tabular}
        \caption{Scores given by each metric for the example summary. $\downarrow$ indicates a lower score is more readable while $\uparrow$ indicates a higher score is more readable. We provide ``S-12'', the score each metric would assign US grade 12, to help contextualize the scores. We additionally translate each score to the US grade level.}
        \label{subtab:ex-metric-scores}
    \end{subtable} 
    \hfill
    \begin{subtable}[t]{.32\textwidth}
     \centering \footnotesize
        \begin{tabular}{r|l}
        \multicolumn{2}{c}{} \\
        \multicolumn{2}{c}{} \\
        \multicolumn{2}{c}{} \\
        Model & Score \\ \hline
        Mistral 7B & 4 \\
        Mixtral 7B & 4.5 \\
        Gemma 7B & 4 \\
        Llama 3.1 8B & 4 \\
        Llama 3.3 70B & 4 \\
        \end{tabular}
    \caption{Scores given by each model for the example summary %\daniel{from X dataset}. 
    The scores are on a scale of 1-5, with 5 being the most readable.}
    \label{subtab:ex-lm-scores}
    \end{subtable} 
    \caption[Example summary with each metrics' readability score.]{ \ref{subtab:ex-summ} contains an example summary from~\textcite{August2024KnowYA}'s dataset. \ref{subtab:ex-metric-scores} contains each metric's score for the example summary. \ref{subtab:ex-lm-scores} contains each model's readability scoring for the example summary. On average, the human annotators rated this summary a 4.05/5, indicating they found the summary fairly readable. All the LM evaluators rate the summary a 4 or 4.5 out of 5, agreeing with the human annotators. In contrast, 6 out of 8 of the traditional metrics rate the summary at a college reading level or higher, which is considered low readability. 
    % \daniel{Say how this is different from metrics/LLMs judgement?}
    }
    \label{tab:example-summ}
\end{table*}
\clearpage
\begin{figure}[h]
    \centering
    \vspace{-2mm}
    \includegraphics[width=.99\columnwidth]{research-chapters/5-chapter/figures/best-prompt.pdf}
        \caption[The best performing prompt of the 3 we tested.]{The best performing prompt of the 3 we tested. We report the results of this prompt in~\Cref{subtab:lm-corr-results} and the results of the remaining prompts in~\Cref{appendix:prompt}.}
    \label{fig:best-prompt}
\end{figure}

\subsection[LMs as evaluators of readability]{LMs as evaluators of readability \hyperref[rq3]{\rqthree}}\label{subsec:results-rqthree}
Traditional readability metrics rely on lexical proxies and do not measure other elements of a summary that could make it more readable, such as definitions of technical terms, explanations of important concepts, or descriptions of impact and motivation. LMs have been shown to perform well on many language understanding tasks~\cite{Brown2020LanguageMA,srivastava2023beyond}, indicating that they have some understanding of language. We hypothesize that this knowledge will translate well to the task of PLS, and the LMs will be able to reason about more complex features of a summary that impact the readability.


We experiment with 3 prompts. We report the best performing prompt in~\Cref{fig:best-prompt} and its results in~\Cref{subtab:lm-corr-results}; the other prompts and their results are in~\Cref{appendix:prompt}.  All of the LMs outperform the traditional metrics in correlation with human judgments. The best performing model, \texttt{Llama 3.3 70B}, outperforms the best traditional metric, DCRS, by nearly $0.2$ Pearson points. We conduct significance testing and report the p-values comparing the LM results to the traditional metrics in \Cref{appendix:sig-testing}. 

Performance in this task is not solely a factor of model size, as we see that smaller models perform similarly to the larger models. The difference in performance between the LMs is small, indicating that most generally well-performing models can be good judges of readability.

\Cref{tab:example-summ} contains an example summary and its associated scores from each LM. The human annotators rated the example summary a 4.05 out of 5 on reading ease. All models gave the summary a rating of 4 or 4.5  out of 5. The reasoning provided by \texttt{Llama 3.3 70B} states that the ``concepts discussed, such as analyzing breath samples and identifying chemical compounds, are also explained in a way that is easy to understand.''
The model notes that the summary ``may require some effort and attention,'' contributing to the model's reasoning for assigning the summary a 4/5 rather than a 5/5. This output indicates that the model is using its language reasoning abilities to rate the summary on attributes deeper than lexical features.


\begin{figure*}[h]
\centering\includegraphics[width=.99\textwidth]{research-chapters/5-chapter/figures/lm_readability-scores.pdf}
    \vspace{-3mm}
    \caption[Histogram of LM readability scores and the mean scores for each dataset.]{Histogram of LM readability scores and the mean scores ($\mu$) for each dataset, as judged by \texttt{Llama 3.3 70B}. As we can see from the results, PLOS and CELLS are judged to be similarly readable to the expert-targeted datasets (arXiv, PubMed, and SciTLDR). The most readable PLS datasets are CDSR and SciNews.}
    \vspace{-4mm}
    \label{fig:histogram-datasets}
\end{figure*}


\begin{table}[t]
\centering \footnotesize
\begin{tabular}{r|ccc}
Dataset & Mean & Median & Var \\ \hline
arXiv & 1.31 & 1 & 0.23 \\
PubMed & 1.99 & 2 & 0.19 \\
SciTLDR & 1.86 & 2 & 0.32 \\
SKJ & 4.40 & 4 & 0.24 \\ 
CDSR & 3.49 & 4 & 0.52 \\
PLOS & 2.06 & 2 & 0.26 \\
eLife & 3.18 & 3 & 0.65 \\
Eureka & 3.21 & 3 & 0.67 \\
CELLS & 2.23 & 2 & 0.50 \\
SciNews & 3.37 & 4 & 0.64
\end{tabular}
\vspace{-2mm}
\caption[LM readability scores on a scale of 1 to 5.]{Readability scores on a scale of 1 to 5, as judged by \texttt{Llama-3.3-70B}, 5 being the most readable. We report the mean, median, and variance of each score.}\label{tab:lm-readability-scores-by-dataset}
\vspace{-4mm}
\end{table}
\subsection[Analysis of summarization datasets]{Analysis of summarization datasets \hyperref[rq4]{\rqfour}}\label{subsec:results-rqfour}
We analyze scientific summarization datasets using the LM evaluators. We use \texttt{Llama 3.3 70B}, the best performing model from \Cref{subsec:results-rqtwo}. In \Cref{fig:histogram-datasets}, we include histograms of the readability scores for all 10 tested datasets, to visualize the distributions. 
In~\Cref{tab:lm-readability-scores-by-dataset}, we report the mean, median, and variance of the readability scores for each dataset.


% In \Cref{appendix:additional-stats}, we report the mean, median, and variance of the LM readability scores. 


We've shown that most LM judgments of readability correlate higher with human judgments than traditional metrics. In order to further validate our findings, we begin our analysis with 4 datasets with specific target audiences - experts or kids. 

\paragraph{Expert-targeted datasets.} We experiment with 3 datasets intended for expert readers: arXiv, PubMed, and SciTLDR. ArXiv, PubMed, and SciTLDR receive low readability scores, averaging less than 2/5. This matches our expectations since summaries intended for an expert audience typically have low readability for non-experts. We also note that SciTLDR receives similarly low readability scores, despite containing significantly shorter summaries than the arXiv and PubMed datasets. This shows that the LM evaluator is not simply favoring shorter summaries as more readable. 

\paragraph{Kid-targeted dataset.} SJK receives high readability scores, with an average readability of $4.40$. The results of the expert and kid targeted datasets match our expectations of readability scores, and serve to support the analysis of the remaining 6 general-audience datasets below.

\paragraph{General audience datasets.} We analyze 6 popular PLS datasets: CDSR, PLOS, eLife, Eureka, CELLS, and SciNews. PLOS and CELLS receive mean readability scores of 2.06 and 2.23, respectively. These scores are similar to the expert-targeted datasets described above, indicating that these two datasets may not be well-suited for PLS. SciNews and CDSR receive the highest readability scores, with average scores of 3.49 and 3.37, respectively, indicating that they are well suited for the task of PLS. 


\begin{figure}[!h]
    \centering
    % First subfigure
    \begin{subfigure}[b]{0.45\textwidth}
        \centering
        \includegraphics[width=\textwidth]{research-chapters/5-chapter/figures/Score_Keyword_Analysis.pdf}
        \vspace{-6mm}
        \caption{Keywords stratified by score.}
        \label{fig:keywords-score}
    \end{subfigure}
    \hfill
    % Second subfigure
    \begin{subfigure}[b]{0.45\textwidth}
        \centering
        \includegraphics[width=\textwidth]{research-chapters/5-chapter/figures/Dataset_Keyword_Analysis.pdf}
        \vspace{-6mm}
        \caption{Keywords stratified by dataset.}
        \label{fig:keywords-dataset}
    \end{subfigure}

    \caption[Keywords mentioned in the reasoning of the LM evaluator for why a summary was given a certain readability score.]{Keywords mentioned in the reasoning of the LM evaluator for why a summary was given a certain readability score. \Cref{fig:keywords-score} contains the keywords stratified by score and \Cref{fig:keywords-dataset} contains keywords stratified by dataset.}
    \label{fig:keyword-analysis}
    \vspace{-6mm}
\end{figure}



\paragraph{Keyword analysis.} To understand the LM's reasoning for assigning scores, we use the \texttt{YAKE!} algorithm to extract keywords from the reasoning provided by the LM evaluator for why each summary was provided with a specific score~\cite{CAMPOS2020257}. \Cref{fig:keywords-score} contains the keywords stratified by score and \Cref{fig:keywords-dataset} contains the keywords stratified by dataset. 
When stratified by score, the model mentions issues such as requiring ``expert background knowledge'' and ``using specialized terms'' for summaries with readability scores of 1 or 2. For summaries with scores of 4 or 5, the model references how the summaries include ``explanations for the non-expert'' and explains ``concepts in an accessible'' manner.
When stratified by dataset, for datasets with generally low readability scores, the model mentions issues such as requiring ``specialized knowledge'' or that the text is ``likely for an academic.'' The model also mentions the domain specific knowledge required such as Pubmed's focus on ``disease[s].'' For datasets with generally high readability scores, such as SJK and SciNews, the model mentions how the summaries are ``easy for most adults'' and how the text is ``understandable by those without'' background knowledge. This keyword analysis indicates LMs are attributing their judgments to deeper attributes that contribute to readability compared to traditional metrics.  

% We report the extracted keywords and conduct a similar analysis, but with the keyword extraction stratified by dataset, in \cref{appendix:keyword-analysis-dataset}. LMs are attributing their judgments to deeper attributes that contribute to readability compared to traditional metrics. 

\begin{table}[h]
\vspace{-2mm}
\centering \footnotesize
\begin{tabular}{r|ll|lll}
 & \multicolumn{2}{c|}{LM Eval} & \multicolumn{3}{c}{FKGL} \\
Dataset & S & R & S & R & $\Delta$R \\ \hline
arXiv & 1.31 & 10 & 11.53 & 2 & $+8$ \\
PubMed & 1.99 & 8 & 14.14 & 5 & $+3$ \\
SciTLDR & 1.86 & 9 & 15.66 & 10 & $-1$ \\
SKJ & 4.40 & 1 & 8.41 & 1 & $0$ \\
CDSR & 3.49 & 2 & 14.08 & 4 & $-2$ \\
PLOS & 2.06 & 7 & 15.44 & 9 & $-2$ \\
eLife & 3.18 & 5 & 11.87 & 3 & $+2$ \\
Eureka & 3.21 & 4 & 14.87 & 6 & $-2$ \\
CELLS & 2.23 & 6 & 15.35 & 8 & $-2$ \\
SciNews & 3.37 & 3 & 14.98 & 7 & $-4$
\end{tabular}
\caption[The mean score and rank for each dataset, as judged by an LM evaluator and FKGL.]{The mean score (S) and rank (R) for each dataset, as judged by an LM evaluator and FKGL. $\Delta$R represents the change in rank from the LM evaluator to the FKGL scores.}\label{tab:ranked-datasets}
\end{table}

\paragraph{LM evaluators vs. traditional metrics.} Finally, we compare the results of the analysis using traditional metrics to LM evaluators of readability. In this analysis, we focus on \texttt{Llama 3.3 70B}, the best performing LM, and FKGL, the most popular readability metric. \Cref{tab:ranked-datasets} compares the average LM readability and FKGL score for each dataset, and how each metric would rank the datasets. All but 1 dataset changed its ranking depending on the metric used. arXiv has the largest delta, ranking 10th in readability according to the LM evaluator and 2nd according to FKGL. FKGL ranking arXiv as the 2nd most readable is particularly concerning, as this dataset is a collection of scientific abstracts, intended for an expert audience.
To measure disagreement, we convert each metric into binary scores of ``high readability'' and ``low readability.'' For FKGL, we consider any summary given a score of under 12 points to have high readability. FKGL considers any text above 12 to be college reading level. For the LM evaluator, we consider any summary given a score of 3 or higher to have high readability. By converting the scores to binary labels, we calculate the Cohen's Kappa score~\cite{McHugh2012InterraterRT} for agreement as 0.17, indicating the two metrics have fair but not substantial agreement. We provide examples of this disagreement in \Cref{tab:disagreement-example}. This analysis shows how the evaluation metrics we use can greatly influence the conclusions we draw. 


\begin{table}[!t]
    \centering
    \begin{subtable}[t]{.44\textwidth}
        \footnotesize
        \begin{tabular}{p{.99\textwidth}}
        \textexample{Wind power is an important source of renewable energy, but some people are concerned that conventional wind turbines are too loud and too hazardous for birds and bats. We wanted to create a new kind of wind energy harvesting machine based on the jiggling motion of cottonwood tree leaves in the wind, which would be quieter and safer for wildlife. After building and testing artificial cottonwood leaves that moved and created electricity in the wind, we found that they didn’t produce enough energy to feasibly use for electricity production. We also tried building a cattail-like device to generate electricity when it swayed in the wind, [...]
        % but it also didn’t produce enough energy to make it reasonable to use. 
        % Though our research showed that artificial plants’ jiggling or swaying isn’t likely to be a cost-effective way to produce electricity, we think it could be fruitful to look into other plant-inspired designs for harvesting wind energy. We also are testing a previously unexploited biological material known to convert mechanical to electrical energy far more effectively than the ones used today.
        } 
        \end{tabular}
        \caption{FKGL = 16.47 (College-graduate), LM score = 4/5.}
        \label{subtab:ex-disaggreement-sjk}
    \end{subtable}%
    \hfill
    \begin{subtable}[t]{.51\textwidth}
        \footnotesize
        \begin{tabular}{p{.99\textwidth}}
        \textexample{Introduction. Accumulation of glycochenodeoxycholic acid (GCDC) in serum has a clinical significance as an inductor of pathological hepatocyte apoptosis, which impairs liver function. Inhibition of GCDC accumulation can be used as a marker in therapy. This study was aimed to quantify the serum level of GCDC in obstructive jaundice patients. Methodology. GCDC acid level in the serum was quantified using high performance liquid chromatography (HPLC) technique according to Muraca and Ghoos modified method. It was performed before and after decompression at day 7 and day 14. The sample was extracted with solid phase extraction (SPE) technique on SPE column. The results were analyzed using SPSS V 16.0 (P < 0.05) [...]
        % and quantified with standard curve on GCDC acid. 
        % Result. There were 21 cases with range of GCDC acid serum level before decompression was 90.9 (SD 205.5) μmol/L and day 7 after decompression decreased to 4.0 (SD 46.4) μmol/L and then increased to 11.3 (SD 21.9) μmol/L (P < 0.05).
        % This method could separate GCDC acid on serum with good resolution, high precision and accuracy, and linear calibration curve on measured level range. Conclusion. HPLC can quantify GCDC acid serum on obstructive jaundice patients and can be used to support its pharmacokinetic study.
        }
        \end{tabular}
        \caption{FKGL = 10.0 (10th grade), \\LM score = 1/5.}
        \label{subtab:ex-disaggreement-plos}
    \end{subtable} 
    \caption[Examples of disagreement between FKGL and the LM evaluator.]{Examples of disagreement between FKGL and the LM evaluator. \ref{subtab:ex-disaggreement-sjk} contains an example from the SJK dataset that the LM rated high readability and FKGL rated low readability. \ref{subtab:ex-disaggreement-plos} contains a summary from the Pubmed dataset that the LM rated low readability while FKGL rated high readability.}
    \label{tab:disagreement-example}
\end{table}

\subsection{Statistical Significance}\label{appendix:sig-testing}
We use the Williams test to calculate statistical significance of the difference in performance between each LM evaluator and traditional metric~\cite{graham-baldwin-2014-testing}. We report the p-values in \Cref{tab:significance-testing}. The difference in Pearson correlation between \texttt{Llama 3.3 70B}, the best performing model, and the traditional metrics is statistically significant, except for DCRS and CLI. The Pearson correlation difference between the LM evaluators and FKGL, the most popular metric, is statistically significant, except \texttt{Llama 3.1 8B}. The Kendall-Tau values show that the Mistral, Mixtral, and Gemma models are statistically significant over most of the traditional metrics. This supports our suggestions from \Cref{subsec:recommendations}, in which we recommend using a combination of the best performing traditional metrics (DCRS and CLI) with LM evaluators, while discontinuing the use of FKGL. 

\begin{table*}[!htb]
\centering
    \begin{subtable}{.9\textwidth}
        \centering \footnotesize
        \begin{tabular}{l|rrrrrrrr}
         & \multicolumn{1}{l}{LW} & \multicolumn{1}{l}{Spache} & \multicolumn{1}{l}{FRE} & \multicolumn{1}{l}{ARI} & \multicolumn{1}{l}{GFI} & \multicolumn{1}{l}{DCRS} & \multicolumn{1}{l}{CLI} & \multicolumn{1}{l}{FKGL} \\ \hline
         Mistral 7B & \cellcolor[HTML]{B7E1CD}6.51E-04 & \cellcolor[HTML]{B7E1CD}0.02 & \cellcolor[HTML]{B7E1CD}0.05 & \cellcolor[HTML]{B7E1CD}0.01 & \cellcolor[HTML]{B7E1CD}0.05 & 0.19 & 0.18 & \cellcolor[HTML]{B7E1CD}0.02 \\
        Mixtral 7B & \cellcolor[HTML]{B7E1CD}6.90E-04 & \cellcolor[HTML]{B7E1CD}0.01 & \cellcolor[HTML]{B7E1CD}0.03 & \cellcolor[HTML]{B7E1CD}0.01 & \cellcolor[HTML]{B7E1CD}0.04 & 0.16 & 0.15 & \cellcolor[HTML]{B7E1CD}0.02 \\
        Gemma 7B & \cellcolor[HTML]{B7E1CD}9.65E-04 & \cellcolor[HTML]{B7E1CD}0.02 & \cellcolor[HTML]{B7E1CD}0.03 & \cellcolor[HTML]{B7E1CD}0.01 & \cellcolor[HTML]{B7E1CD}0.04 & 0.17 & 0.15 & \cellcolor[HTML]{B7E1CD}0.02 \\
        Llama 3.1 8B & \cellcolor[HTML]{B7E1CD}2.00E-03 & \cellcolor[HTML]{B7E1CD}0.04 & 0.14 & \cellcolor[HTML]{B7E1CD}0.03 & 0.10 & 0.34 & 0.31 & 0.06 \\
        Llama 3.1 70B & \cellcolor[HTML]{B7E1CD}3.66E-04 & \cellcolor[HTML]{B7E1CD}0.01 & \cellcolor[HTML]{B7E1CD}0.02 & \cellcolor[HTML]{B7E1CD}0.01 & \cellcolor[HTML]{B7E1CD}0.03 & 0.14 & 0.13 & \cellcolor[HTML]{B7E1CD}0.02 \\
        \end{tabular}
        \caption{Pearson correlation p-values.}\label{subtab:pearson-pval}
    \end{subtable}%
    
    \begin{subtable}[t]{.9\textwidth}
        \centering \footnotesize
          \begin{tabular}{l|rrrrrrrr}
         & \multicolumn{1}{l}{LW} & \multicolumn{1}{l}{Spache} & \multicolumn{1}{l}{FRE} & \multicolumn{1}{l}{ARI} & \multicolumn{1}{l}{GFI} & \multicolumn{1}{l}{DCRS} & \multicolumn{1}{l}{CLI} & \multicolumn{1}{l}{FKGL} \\ \hline
        Mistral 7B & \cellcolor[HTML]{B7E1CD}0.01 & \cellcolor[HTML]{B7E1CD}0.01 & \cellcolor[HTML]{B7E1CD}0.03 & \cellcolor[HTML]{B7E1CD}0.01 & \cellcolor[HTML]{B7E1CD}0.04 & 0.13 & 0.10 & \cellcolor[HTML]{B7E1CD}0.03 \\
        Mixtral 7B & \cellcolor[HTML]{B7E1CD}0.01 & \cellcolor[HTML]{B7E1CD}0.02 & \cellcolor[HTML]{B7E1CD}0.05 & \cellcolor[HTML]{B7E1CD}0.02 & 0.06 & 0.19 & 0.14 & \cellcolor[HTML]{B7E1CD}0.04 \\
        Gemma 7B & \cellcolor[HTML]{B7E1CD}0.01 & \cellcolor[HTML]{B7E1CD}0.02 & \cellcolor[HTML]{B7E1CD}0.03 & \cellcolor[HTML]{B7E1CD}0.02 & \cellcolor[HTML]{B7E1CD}0.05 & 0.16 & 0.10 & \cellcolor[HTML]{B7E1CD}0.03 \\
        Llama 3.1 8B & \cellcolor[HTML]{B7E1CD}0.02 & 0.05 & 0.12 & \cellcolor[HTML]{B7E1CD}0.05 & 0.12 & 0.31 & 0.25 & 0.09 \\
        Llama 3.1 70B & \cellcolor[HTML]{B7E1CD}0.03 & 0.06 & 0.09 & \cellcolor[HTML]{B7E1CD}0.05 & 0.12 & 0.30 & 0.24 & 0.09 \\
        \end{tabular}
        \caption{Kendall-Tau p-values.}\label{subtab:kt-pval}
    \end{subtable} 
    \caption[Williams test p-values comparing the difference in performance between each LM and each traditional metric.]{Williams test p-values comparing the difference in performance between each LM and each traditional metric. Values that are statistically significant ($\textrm{p-value} < 0.05$), are highlighted in green.}
    \label{tab:significance-testing}
\end{table*}